% 张三慧《大学物理学》第4版第二卷 第30章《核物理》
% 主教材PDF物理页485--487（印刷468--470），配套S PDF392--399。
% 仅原书正式章末习题30.1--30.22；思考题30.1--30.10不混入。
\section*{第30章\quad 核物理}
\subsection*{章末习题30.1--30.22}
\paragraph{习题 30.1}
动能为6MeV的\(\alpha\)粒子与初始静止的
\({}^{197}\mathrm{Au}\) 金核正面接近，求最小距离。
改为\({}^{14}\mathrm N\) 氮核时求最小距离；
此时可否忽略核的反冲？此距离是否能进入
\({}^{14}\mathrm N\) 的核力作用范围？
\paragraph{习题 30.2}
\({}^{16}\mathrm N\)、\({}^{16}\mathrm O\) 和
\({}^{16}\mathrm F\) 的\emph{中性原子质量}分别为
16.006099u、15.994915u及16.011465u。
计算三原子所对应原子核的结合能，
注意原子质量及氢原子质量使用的一致性。
\paragraph{习题 30.3}
把核内质子当成三维费米气，推导原子序数
\(Z\)、质量数 \(A\) 原子核内质子的费米能量
与平均动能；对 \({}^{56}\mathrm{Fe}\) 和
\({}^{238}\mathrm U\) 作数值计算（MeV），
取核半径 \(R=r_0A^{1/3}\)，\(r_0=1.2\,\mathrm{fm}\)。
\paragraph{习题 30.4}
三对镜像核
\(({}^{11}\mathrm C,{}^{11}\mathrm B)\)、
\(({}^{15}\mathrm O,{}^{15}\mathrm N)\)、
\(({}^{21}\mathrm{Na},{}^{21}\mathrm{Ne})\)，
静电能差分别为2.79MeV、3.48MeV和4.30MeV。
把核近似为均匀带电球，估计三对核的半径，
是否大致正比 \(A^{1/3}\)，求比例常量。
\paragraph{习题 30.5}
假设 \({}^{12}\mathrm C\) 是三颗\(\alpha\)
在三角形顶点、\({}^{16}\mathrm O\) 是四颗
\(\alpha\) 在四面体顶点形成的团簇核。
若每对\(\alpha\)间结合能为2.42MeV，
并计入\(\alpha\)本身结合能，
与质量亏损计算的两核结合能比较。
原子质量（u）：\({}^{1}\mathrm H=1.007825\)、
\({}^{4}\mathrm{He}=4.002603\)、
\({}^{12}\mathrm C=12.000000\)、
\({}^{16}\mathrm O=15.994915\)。
\paragraph{习题 30.6}
假设 \({}^{232}\mathrm{Th}\) 核分裂成完全相等的
两块（两个 \(A=116,Z=45\) 的核），
用液滴模型半经验核结合能公式计算
释放的能量，并说明表面项、库仑项的贡献。
\paragraph{习题 30.7}
两枚相同质量数 \(A\)、质子数 \(Z\) 的核
聚合成 \(2A,2Z\) 原子核。
由半经验核结合能公式推导反应放能表达式，
求两个 \({}^{12}\mathrm C\) 核聚合时释放多少能量。
\paragraph{习题 30.8}
放射性核的平均寿命为 \(\tau\)。
按指数衰变概率分布求该寿命的方均根偏差，
并求最概然寿命。
\paragraph{习题 30.9}
天然钾中 \({}^{40}\mathrm K\) 丰度
\(1.2\times10^{-4}\)，半衰期
\(1.3\times10^9\)年；钾占人体质量0.37\%。
按一个体重70kg的人估计体内
\({}^{40}\mathrm K\) 的放射性活度。
\paragraph{习题 30.10}
计算10kg铀矿 \(\mathrm U_3\mathrm O_8\) 中
\({}^{226}\mathrm{Ra}\) 与 \({}^{231}\mathrm{Pa}\) 的
质量。天然铀同位素丰度
\({}^{238}\mathrm U=99.27\%\)、
\({}^{235}\mathrm U=0.72\%\)，
\({}^{226}\mathrm{Ra}\) 半衰期1600年，
\({}^{231}\mathrm{Pa}\) 半衰期
\(3.27\times10^4\)年。
可以用铀衰变链长期封闭下的放射性平衡近似；
必要的母核半衰期须明确写出。
\paragraph{习题 30.11}
患者服用活度30\,$\mu$Ci的
\({}^{123}\mathrm I\)，24小时后测得甲状腺
活度4\,$\mu$Ci。已知半衰期13.1小时，
求原服用碘在这24小时内集中于甲状腺的比例。
通常健康人体该比例约15\%--40\%。
\paragraph{习题 30.12}
向患者静脉注射含活度300kBq的
\({}^{24}\mathrm{Na}\) 生理盐水，10小时后血液活度
为每立方厘米30Bq。已知半衰期14.97小时，
求患者全身血液总体积。
\paragraph{习题 30.13}
待测古木片在纯氧中燃烧后收集0.3mol
的CO$_2$，由其中\({}^{14}\mathrm C\) 衰变导致
每分钟计数9次。用天然碳中
\({}^{14}\mathrm C/{}^{12}\mathrm C=1.3\times10^{-12}\)、
\({}^{14}\mathrm C\) 平均寿命8270年，
估算古木年龄（假定探测效率100\%，
未作日历年校正）。
\paragraph{习题 30.14}
某岩石中有0.30g的\({}^{238}\mathrm U\)
和0.12g的\({}^{206}\mathrm{Pb}\)，
假设所有这种铅都由铀衰变产生，
且矿物封闭、无初始铅，
求地质年龄（\({}^{238}\mathrm U\)半衰期
约 \(4.46\times10^9\)年）。
\paragraph{习题 30.15}
\({}^{226}\mathrm{Ra}\) 的\(\alpha\)衰变
发射动能为4.7825MeV的\(\alpha\)粒子，
求子核\({}^{222}\mathrm{Rn}\)反冲动能和该次
\(\alpha\)衰变的总放能 \(Q\)。
\paragraph{习题 30.16}
\({}^{64}\mathrm{Cu}\) 可经
\(\beta^-\) 衰变成\({}^{64}\mathrm{Zn}\)，
或经 \(\beta^+\) 衰变成\({}^{64}\mathrm{Ni}\)。
两衰变的 \(Q\) 值分别为0.57MeV、0.66MeV。
求 \({}^{64}\mathrm{Zn}\) 与
\({}^{64}\mathrm{Ni}\) 的核质量差（u），
注意核质量与中性原子质量约定差别。
\paragraph{习题 30.17}
将天然丰度100\%的
\({}^{59}\mathrm{Co}\) 金属立方块，边长2cm，
置于热中子通量
\(2\times10^{12}\,\mathrm{cm^{-2}s^{-1}}\) 中
照射6小时，使
\({}^{59}\mathrm{Co}+n\to{}^{60}\mathrm{Co}+\gamma\)。
反应截面120b，钴密度8.858g/cm$^3$，
\({}^{60}\mathrm{Co}\) 半衰期5.27年。
估计取出后的活度，并比较薄靶近似与
厚靶中子衰减、自屏蔽限流的物理边界。
\paragraph{习题 30.18}
Cd有8种稳定同位素，平均热中子吸收截面
4000b，摩尔质量112.4g/mol、密度
8.64g/cm$^3$。为吸收95\%入射中子，
求Cd片的厚度；比较
\(1-\exp(-n\sigma d)=0.95\) 的指数衰减结果
与薄靶线性近似。
\paragraph{习题 30.19}
求以下三反应的 \(Q\) 值并判断吸热或放热：
\[
{}^{13}\mathrm C(p,\alpha){}^{10}\mathrm B,\quad
{}^{13}\mathrm C(p,d){}^{12}\mathrm C,\quad
{}^{13}\mathrm C(p,\gamma){}^{14}\mathrm N.
\]
原子质量（u）：
\({}^{13}\mathrm C=13.003355\)、
\({}^{1}\mathrm H=1.007825\)、
\({}^{4}\mathrm{He}=4.002603\)、
\({}^{10}\mathrm B=10.012937\)、
\({}^{2}\mathrm H=2.014102\)、
\({}^{14}\mathrm N=14.003074\)；
\({}^{12}\mathrm C=12.000000\)u。
\paragraph{习题 30.20}
计算反应
\({}^{13}\mathrm C(p,\alpha){}^{10}\mathrm B\)
的实验室入射质子阈能，
除负 \(Q\) 的运动学阈值外，
还应计入质子必须靠近靶核到核半径
\(r=r_0(A^{1/3}+1)\)、\(r_0=1.2\,\mathrm{fm}\)
附近的库仑势垒能量。
质量数据同习题30.19；明确说明两种阈值
相加的简化模型与量子隧穿可能性。
\paragraph{习题 30.21}
太阳内部约含
\(1.5\times10^{30}\,\mathrm{kg}\) 氢，
辐射总功率 \(3.9\times10^{26}\,\mathrm W\)。
在维持此功率及氢均转化为氦的假设下，
估算消耗完氢所需时间。
注明每4个氢核聚变的放能，以及中微子
逃逸、电子正电子湮灭对实际可供辐射能的影响。
\paragraph{习题 30.22}
较太阳热的恒星里通过CNO碳氮氧循环融合氢，
按原书六步：
\[
\begin{aligned}
{}^{1}\mathrm H+{}^{12}\mathrm C&\to{}^{13}\mathrm N+\gamma,\\
{}^{13}\mathrm N&\to{}^{13}\mathrm C+e^++\nu_e,\\
{}^{1}\mathrm H+{}^{13}\mathrm C&\to{}^{14}\mathrm N+\gamma,\\
{}^{1}\mathrm H+{}^{14}\mathrm N&\to{}^{15}\mathrm O+\gamma,\\
{}^{15}\mathrm O&\to{}^{15}\mathrm N+e^++\nu_e,\\
{}^{1}\mathrm H+{}^{15}\mathrm N&\to{}^{12}\mathrm C+{}^{4}\mathrm{He}.
\end{aligned}
\]
(1) 证明碳为催化循环而无净消耗；
(2) 计算每步核反应的 \(Q\) 值，指出
\(\beta^+\) 的中性原子质量应扣除两倍
电子静质量；(3) 计算六步之和及计入
两次正电子与环境电子湮灭后的总核融合放能。
原子质量（u）：\({}^{1}\mathrm H=1.007825\)、
\({}^{13}\mathrm N=13.005738\)、
\({}^{13}\mathrm C=13.003355\)、
\({}^{14}\mathrm N=14.003074\)、
\({}^{15}\mathrm O=15.003065\)、
\({}^{15}\mathrm N=15.000109\)、
\({}^{4}\mathrm{He}=4.002603\)，
\({}^{12}\mathrm C=12\)u。
